\documentclass[11pt]{article}
\usepackage[margin=1in]{geometry}
\usepackage{enumitem}
% =============================================================================
% Preambles/header.tex — shared LaTeX/Beamer preamble for this project
%
% Use from a Beamer lecture by prepending:
%     \input{header}
% and compiling with TEXINPUTS=../Preambles:$TEXINPUTS (the /compile-latex
% skill does this automatically).
%
% PALETTE CONTRACT. Color names defined here MUST match the names in
% Quarto/theme-template.scss so Beamer and Quarto renderings use the same
% palette. The scripts/check-palette-sync.sh script enforces this.
%
% To customize: edit the HEX values below AND the matching $variable in
% Quarto/theme-template.scss, then run scripts/check-palette-sync.sh.
% =============================================================================

% -----------------------------------------------------------------------------
% Engine detection + encoding
%
% Our /compile-latex skill uses XeLaTeX, where inputenc is unnecessary and
% can produce encoding warnings. Only load inputenc under pdfTeX.
% -----------------------------------------------------------------------------
\usepackage{iftex}
\ifPDFTeX
  \usepackage[utf8]{inputenc}
\fi

\usepackage{amsmath, amssymb, amsthm}
\usepackage{booktabs}
\usepackage{graphicx}

% -----------------------------------------------------------------------------
% PALETTE — mirrors Quarto/theme-template.scss variable names.
% Load xcolor and define colors BEFORE anything that references them
% (notably hyperref, hypersetup, and the Beamer color assignments).
% -----------------------------------------------------------------------------
\usepackage{xcolor}

\definecolor{primary-blue}{HTML}{012169}      % headings, accents
\definecolor{primary-gold}{HTML}{B9975B}      % emphasis, borders
\definecolor{highlight-yellow}{HTML}{F2A900}  % markers, alerts
\definecolor{light-bg}{HTML}{E8EDF5}          % title / section backgrounds
\definecolor{jet}{HTML}{1A1A1A}               % body text

% Semantic colors (used in diagrams and annotations)
\definecolor{positive}{HTML}{15803D}          % good / observed / identified
\definecolor{negative}{HTML}{B91C1C}          % bad / problematic / bias
\definecolor{neutral}{HTML}{525252}           % reference / context

% Highlight accents (match .hi-* classes in the SCSS)
\definecolor{hi-slate}{HTML}{314F4F}
\definecolor{hi-green}{HTML}{2E8B57}
\definecolor{hi-red}{HTML}{C41E3A}

% Semantic aliases so skill templates and snippets can write
% \textcolor{accent}{...} without hard-coding "primary-blue".
\colorlet{accent}{primary-blue}
\colorlet{accent2}{primary-gold}

% -----------------------------------------------------------------------------
% Hyperref — loaded AFTER xcolor + palette so link colors resolve.
% -----------------------------------------------------------------------------
\usepackage{hyperref}
\hypersetup{colorlinks=true, linkcolor=primary-blue, urlcolor=primary-blue,
            citecolor=primary-blue, pdfborder={0 0 0}}

% -----------------------------------------------------------------------------
% Course code — set once per deck (after \input{header}, before \begin{document})
% via \coursecode{CS401}. Shown in the footer on every slide so a deck's course
% is identifiable without opening the title page. Empty by default.
% -----------------------------------------------------------------------------
\makeatletter
\newcommand{\coursecode}[1]{\gdef\@coursecodetext{#1}}
\gdef\@coursecodetext{}
\makeatother

% -----------------------------------------------------------------------------
% Beamer theme assignments (only apply under Beamer).
%
% We detect Beamer via \@ifclassloaded{beamer}, which is the documented,
% non-internal way. This must be wrapped in \makeatletter/\makeatother
% because \@ifclassloaded contains an @-catcode character.
% -----------------------------------------------------------------------------
\makeatletter
\@ifclassloaded{beamer}{%
  \usefonttheme{professionalfonts}
  \setbeamercolor{structure}{fg=primary-blue}
  \setbeamercolor{frametitle}{fg=primary-blue}
  \setbeamercolor{title}{fg=primary-blue}
  \setbeamercolor{subtitle}{fg=primary-gold}
  \setbeamercolor{author}{fg=primary-blue}
  \setbeamercolor{institute}{fg=neutral}
  \setbeamercolor{date}{fg=primary-gold}
  \setbeamercolor{itemize item}{fg=primary-blue}
  \setbeamercolor{itemize subitem}{fg=primary-gold}
  \setbeamercolor{alerted text}{fg=primary-gold}
  \setbeamercolor{block title}{fg=white, bg=primary-blue}
  \setbeamercolor{block body}{bg=light-bg}
  % Semantic block variants: block=definition/concept (blue), exampleblock=
  % worked example (green/positive), alertblock=key takeaway or caution (gold).
  % Keep to <=2 colored boxes per slide across ALL three variants (INV-7).
  \setbeamercolor{block title example}{fg=white, bg=positive}
  \setbeamercolor{block body example}{bg=light-bg}
  \setbeamercolor{block title alerted}{fg=white, bg=primary-gold}
  \setbeamercolor{block body alerted}{bg=light-bg}

  % Minimal footer: course code (left) + page X / N (right), no navigation clutter
  \setbeamertemplate{navigation symbols}{}
  \setbeamertemplate{footline}{%
    \color{neutral}\footnotesize\hspace{0.7em}\@coursecodetext\hfill
    \insertframenumber/\inserttotalframenumber\hspace{0.7em}\vspace{0.3em}}
}{}
\makeatother

% -----------------------------------------------------------------------------
% TikZ setup — libraries the snippet gallery uses
% -----------------------------------------------------------------------------
\usepackage{tikz}
\usetikzlibrary{arrows.meta, positioning, calc, decorations.pathreplacing,
                fit, shapes.geometric, backgrounds}

% Shared TikZ styles — snippets and hand-written diagrams can pick these up
% via `\begin{tikzpicture}[every node/.style={...}]` or by naming specific
% styles (e.g., `\node[dag-node] (X) at (0,0) {X};`).
\tikzset{
  % Node styles for causal DAGs and similar
  dag-node/.style = {
    draw=primary-blue, fill=light-bg, rounded corners=2pt,
    minimum width=1.2cm, minimum height=0.9cm, align=center, font=\small
  },
  decision-node/.style = {
    draw=primary-gold, fill=white, diamond, aspect=1.8,
    minimum width=1.8cm, minimum height=1.2cm, align=center, font=\small,
    inner sep=1pt
  },
  % Edge styles — observed vs counterfactual
  observed-edge/.style = {-{Latex[length=2.5mm]}, thick, draw=jet},
  counterfactual-edge/.style = {-{Latex[length=2.5mm]}, thick, draw=neutral, dashed},
  confound-edge/.style = {-{Latex[length=2.5mm]}, thick, draw=negative, dashed},
  % Data-point styles for plots
  observed-dot/.style = {fill=primary-blue, draw=primary-blue, circle, inner sep=1.2pt},
  counterfactual-dot/.style = {fill=white, draw=neutral, circle, inner sep=1.2pt}
}

% -----------------------------------------------------------------------------
% Convenience macros
% -----------------------------------------------------------------------------
\newcommand{\muted}[1]{\textcolor{neutral}{#1}}
\newcommand{\key}[1]{\textcolor{primary-gold}{\textbf{#1}}}
\newcommand{\good}[1]{\textcolor{positive}{#1}}
\newcommand{\bad}[1]{\textcolor{negative}{#1}}

% Standout frame macro: full-bleed dark background for section transitions.
% Beamer-only (uses \begin{frame}/\setbeamercolor/\usebeamerfont) -- do not
% call from an article-class file (e.g. Notes/<CODE>/*-notes.tex). \muted,
% \key, \good, \bad, and \coursecode above are plain xcolor/gdef and are
% safe to use from either class; verified when Notes/article-class support
% was added.
% Expand: \transitionslide{Your transition title}
\newcommand{\transitionslide}[1]{%
  \begingroup
  \setbeamercolor{background canvas}{bg=primary-blue}
  \begin{frame}[plain]
    \centering
    \vfill
    {\usebeamerfont{title}\color{white}\Large #1\par}
    \vfill
  \end{frame}
  \endgroup
}

% Section divider frame (Beamer-only), an alternative to \transitionslide with a
% darker two-line style: near-black (jet) background, a small white label line
% above a larger gold title, and a thin gold rule along the bottom edge.
% Expand: \sectiondivider{Section 2}{Pointers}
\newcommand{\sectiondivider}[2]{%
  \begingroup
  \setbeamercolor{background canvas}{bg=jet}
  \begin{frame}[plain]
    \vspace*{3.0cm}
    {\color{white}\ttfamily\large #1\par}
    \vspace{0.5cm}
    {\color{primary-gold}\LARGE #2\par}
    \begin{tikzpicture}[remember picture, overlay]
      \draw[primary-gold, line width=2.5pt]
        ($(current page.south west)+(0,0.1)$) -- ($(current page.south east)+(0,0.1)$);
    \end{tikzpicture}
  \end{frame}
  \endgroup
}

% -----------------------------------------------------------------------------
% Channel watermark — "Concepts That Click" (YouTube).
%
% Article-class files (CompetitiveExam Q/A sets, Notes, Assignments) get a
% light diagonal draft watermark on every page plus a centered footer naming
% the channel. Beamer decks get a subtle TikZ background watermark instead
% (the footline already carries the course code). Centralized here so every
% MCQ PDF is branded without per-file edits.
% -----------------------------------------------------------------------------
\makeatletter
\@ifclassloaded{article}{%
  \usepackage{draftwatermark}
  \SetWatermarkText{Concepts That Click}
  \SetWatermarkScale{1}
  \SetWatermarkFontSize{2cm}
  \SetWatermarkLightness{0.86}
  \SetWatermarkAngle{45}
  \usepackage{fancyhdr}
  \pagestyle{fancy}
  \fancyhf{}
  \fancyfoot[C]{\footnotesize\textcolor{neutral}{Concepts That Click~|~YouTube: \href{https://www.youtube.com/@concepts.thatclick}{@concepts.thatclick}}}
  \renewcommand{\headrulewidth}{0pt}
  \renewcommand{\footrulewidth}{0pt}
  \fancypagestyle{plain}{%
    \fancyhf{}%
    \fancyfoot[C]{\footnotesize\textcolor{neutral}{Concepts That Click~|~YouTube: \href{https://www.youtube.com/@concepts.thatclick}{@concepts.thatclick}}}%
    \renewcommand{\headrulewidth}{0pt}%
    \renewcommand{\footrulewidth}{0pt}%
  }
}{}
\@ifclassloaded{beamer}{%
  \setbeamertemplate{background}{%
    \begin{tikzpicture}[remember picture, overlay]
      \node[rotate=45, opacity=0.055, align=center]
        at (current page.center)
        {\Huge\textcolor{neutral}{Concepts That Click}};
    \end{tikzpicture}%
  }
}{}
\makeatother 
\coursecode{JKSSB}
\renewcommand{\thesection}{3.\arabic{section}}
\newtheorem{example}{Example}[section]
\newenvironment{definitionbox}[1]{%
  \par\noindent\textbf{Definition (#1).}\ \itshape
}{\par\vspace{0.3em}}
\title{Bits, Bytes \& Units --- Detailed Lecture Notes}
\author{JKSSB Computer Awareness}
\date{\today}

\begin{document}
\maketitle
\noindent\textbf{Video lesson:}
\href{https://youtu.be/j_62CGCBElo}{Bits, Bytes \& KB--TB Explained}

\section{The basic units of digital data}
Digital information is represented using bits. Groups of bits form larger
units, but the terms do not all describe the same kind of quantity. A bit is
a single binary digit; a word is an architecture-dependent processor unit;
a character is an encoded symbol.

\begin{definitionbox}{Bit}
A binary digit that has one of two values, 0 or 1; the smallest unit of data
in the usual computer-fundamentals description.
\end{definitionbox}

\begin{definitionbox}{Nibble}
A group of four bits.
\end{definitionbox}

\begin{definitionbox}{Byte}
A group of eight bits, equal to two nibbles.
\end{definitionbox}

\begin{definitionbox}{Word}
The natural data size a processor handles as a unit; its size depends on the
processor architecture.
\end{definitionbox}

\begin{definitionbox}{Character}
A text symbol represented in digital storage by a character encoding; its
storage size is not necessarily one byte.
\end{definitionbox}

\begin{center}
\begin{tabular}{p{0.19\textwidth}p{0.2\textwidth}p{0.48\textwidth}}
\toprule
\textbf{Unit} & \textbf{Size / basis} & \textbf{Meaning}\\
\midrule
Bit & 1 binary digit & One of the two values 0 or 1\\
Nibble & 4 bits & Half of an 8-bit byte\\
Byte & 8 bits & Two nibbles; common unit for file/storage amounts\\
Word & Architecture-dependent & Natural processor data size, such as 32 or 64 bits\\
Character & Encoding-dependent & A text symbol; its encoded storage may take more than one byte\\
\bottomrule
\end{tabular}
\end{center}

The bit-to-byte relationships are fixed in this lesson:
\[
  1\text{ nibble}=4\text{ bits},\qquad
  1\text{ byte}=8\text{ bits}=2\text{ nibbles}.
\]
The word size is not defined by this ladder. A processor may work with a
natural word size such as 32 or 64 bits, but the term depends on the
architecture and should not be treated as one universal number.

\begin{example}
Write the byte \texttt{10110110} as two nibbles: \texttt{1011} and
\texttt{0110}. Each group contains four bits; together the two groups still
contain eight bits, so the original value is one byte.
\end{example}

\section{Bits and bytes in names and data rates}
Capitalization changes the unit. Lowercase \texttt{b} denotes a bit, while
uppercase \texttt{B} denotes a byte. A byte contains eight bits. Therefore a
rate reported in megabits per second is not numerically equal to a rate in
megabytes per second.

\begin{definitionbox}{Transfer rate}
The amount of data transferred in a stated time interval, commonly written
in bits per second or bytes per second.
\end{definitionbox}

\begin{definitionbox}{Capacity}
The amount of data that a storage medium can hold.
\end{definitionbox}

The slash in a rate such as Mbps means ``per second.'' It is not part of the
data amount itself. Keep the numerical amount, whether it is bits or bytes,
and the time unit visible while calculating. Under matching decimal prefixes,
\[
  1\text{ byte}=8\text{ bits},\qquad
  1\text{ MB}=8\text{ Mb}.
\]
Thus dividing a rate in Mb/s by eight gives MB/s, while multiplying a file
size in MB by eight gives Mb. These conversions change the unit, not the
underlying amount of data.

To convert a bit rate to a byte rate, divide by eight, provided the prefixes
match and the question is asking for the ideal conversion. Protocol overhead,
congestion, and other factors can reduce actual network throughput, so a
nominal rate is not a promise of application-level speed.

\begin{example}
A link is rated at $8$ Mbps. Using matching decimal prefixes and ignoring
overhead, convert bits to bytes:
\[
8\text{ megabits/s}\div 8=1\text{ megabyte/s}.
\]
This is a transfer rate, not a storage capacity.
\end{example}

\begin{example}
Convert an ideal rate of $24$ Mbps to megabytes per second. The prefixes
already match, so only the bit-to-byte conversion is needed:
\[
  24\text{ Mb/s}\div 8=3\text{ MB/s}.
\]
The answer is a rate because the time denominator remains. If the question
asked how long a file takes to transfer, the file amount would first need to
be expressed in the same unit as the rate, then divided by that rate.
\end{example}

\begin{example}
A $4$ MB file contains $32$ megabits under matching decimal prefixes because
$4\text{ MB}\times 8=32\text{ Mb}$. If transferred over an ideal $20$ Mbps
link, the simple lower-bound calculation is $32/20=1.6$ seconds. Actual
transfer can take longer because the ideal calculation ignores protocol
overhead and other delays.
\end{example}

\section{Capacity units and the 1024 convention}
For the conventional JKSSB-style ladder used in this lesson, each step is a
factor of 1024:
\[
1\text{ KB}=1024\text{ bytes},\quad
1\text{ MB}=1024\text{ KB},\quad
1\text{ GB}=1024\text{ MB},\quad
1\text{ TB}=1024\text{ GB}.
\]
Thus $1$ MB under this convention is $1024\times1024=1{,}048{,}576$ bytes.
The exam-style abbreviations KB, MB, GB, and TB are often used with this
traditional binary ladder. Modern technical standards distinguish decimal
SI units (powers of 1000) from binary IEC units such as KiB (powers of 1024).
If a question specifies decimal units or the IEC prefixes, use those
definitions instead. Do not silently mix conventions.

\begin{center}
\begin{tabular}{p{0.18\textwidth}p{0.25\textwidth}p{0.43\textwidth}}
\toprule
\textbf{Unit} & \textbf{Previous unit} & \textbf{Bytes under repeated 1024 convention}\\
\midrule
1 KB & 1024 bytes & $1{,}024$ bytes\\
1 MB & 1024 KB & $1{,}048{,}576$ bytes\\
1 GB & 1024 MB & $1{,}073{,}741{,}824$ bytes\\
1 TB & 1024 GB & $1{,}099{,}511{,}627{,}776$ bytes\\
\bottomrule
\end{tabular}
\end{center}

To move from a larger unit to the next smaller unit, multiply by 1024.
Repeat once for every step crossed. To move from a smaller unit to the next
larger unit, divide by 1024. Writing the units beside the numbers helps
prevent using the wrong direction.

\begin{example}
Convert $3$ KB to bytes under the 1024 convention:
\[
3\text{ KB}\times 1024\text{ bytes/KB}=3072\text{ bytes}.
\]
The KB units cancel, leaving bytes.
\end{example}

\begin{example}
Convert $2$ MB to KB under the same convention:
\[
2\text{ MB}\times 1024\text{ KB/MB}=2048\text{ KB}.
\]
This is one step down the capacity ladder, so multiply by 1024.
\end{example}

\begin{example}
Convert $5$ GB to MB using the same convention. GB is one step larger than
MB, so multiply by 1024:
\[
  5\text{ GB}\times 1024\text{ MB/GB}=5120\text{ MB}.
\]
The units GB cancel, leaving MB. To convert $5120$ MB back to GB, divide by
1024 and recover $5$ GB.
\end{example}

\section{Keep capacity separate from rate}
Capacity answers ``how much can be stored?'' Transfer rate answers ``how
much is moved per unit of time?'' MB is a byte-based amount; Mbps is a
bit-based rate. The lowercase/uppercase distinction and the \texttt{/s}
time component carry information. A drive with a larger capacity does not
necessarily transfer data faster than a smaller drive.

\begin{center}
\begin{tabular}{p{0.2\textwidth}p{0.31\textwidth}p{0.31\textwidth}}
\toprule
\textbf{Question} & \textbf{Capacity} & \textbf{Transfer rate}\\
\midrule
What does it measure? & Amount that can be held & Amount moved per unit of time\\
Common example & A 4 MB file or 1 TB drive & A 20 Mbps connection\\
Unit clue & No time denominator & Includes ``per second'' or \texttt{/s}\\
\bottomrule
\end{tabular}
\end{center}

\begin{example}
In the lesson's opening comparison, $4$ MB describes the file's storage
amount, while $20$ Mbps describes the network's transfer rate. The first
has no time denominator; the second is measured per second, so they are
different kinds of quantities.
\end{example}

Before calculating, perform a quick unit check: if the final question asks
for bytes, convert bits to bytes; if it asks for seconds, divide an amount
of data by a compatible rate; if it asks for a storage amount, do not attach
a time unit. A numerically correct calculation with mismatched units is not
a correct answer.

\subsection*{Exercises}
\begin{enumerate}[label=\arabic*.]
  \item How many bits are in $7$ bytes? Show the conversion.
  \item How many bytes are in $5$ KB under the 1024 convention?
  \item Convert $24$ Mbps to MB/s using matching decimal prefixes and
    ignoring overhead.
  \item A device stores $4$ GB. Is this a capacity or a transfer rate? What
    unit clue supports your answer?
  \item A processor has a 64-bit word size. How many bytes are in one word?
\end{enumerate}

\subsection*{Solutions}
\begingroup\small
\begin{enumerate}[label=\arabic*.]
  \item $7\times 8=56$ bits.
  \item $5\times 1024=5120$ bytes.
  \item $24/8=3$ MB/s before protocol overhead.
  \item It is capacity: GB is a data amount, and there is no ``per second''
    time component.
  \item $64/8=8$ bytes per word.
\end{enumerate}
\endgroup
\end{document}
